% !PRINTED 454
\weberpage{454}
%<<CONT
there remain only the following two possibilities for the fulfilment of (23):

1.\quad $D=-3, \ x=1, \ y=\pm 1.$

Since for the discriminant $-3$ there is only one class of forms, we may assume $A=1$, $B=1$, $C=1$, i.e. we may set for $\omega$ the imaginary third root of unity $e^{\frac{2\pi i}{3}}$, and we obtain among the quantities (15) three equivalent ones by assuming $y=+1$, and $=-1$: % [sic] German print omits y before the second equality

\weq{24}{p\,\omega, \quad \frac{\omega}{p}, \quad \frac{\omega+1}{p},}

as is also recognised from the equations:
\[
\frac{\omega}{p} = \frac{-1}{p+p\omega}, \quad \frac{\omega+1}{p} = \frac{-1}{p\omega}
\]
from which the equivalence of the three quantities (24) is evident.

2.\quad $D=-4, \ x=1, \ y=0.$

We may likewise assume $B=0$, $A=C=1$, i.e. $\omega=i$, and we find the two equivalent values:

\weq{25}{p\,\omega, \quad \frac{\omega}{p},}

as also follows from the equation:
\[
p\,\omega = \frac{-p}{\omega}
\]

Let, further, two of the values (3) % [sic] German print has (3); the quantities under discussion are (15)
\[
\frac{\omega+c}{p}, \quad \frac{\omega+c'}{p}
\]
be equivalent, so that an equation holds:

\weq{26}{\frac{\omega+c}{p} = \frac{\gamma p+\delta(\omega+c')}{\alpha p+\beta(\omega+c')}, \quad \alpha\delta-\beta\gamma=1,}

or

\weq{27}{\begin{aligned} \beta\,\omega^2 + (\beta c'+\beta c+\alpha p-\delta p)\omega \; +\phantom{0}\\ \beta c c'+\alpha c p-\delta c' p-\gamma p^2 &= 0, \end{aligned}}

from which, by comparison with (9),

\weq{28}{\begin{aligned} \beta &= Ax,\\ \beta c'+\beta c+\alpha p-\delta p &= Bx,\\ \beta c c'+\alpha c p-\delta c' p-\gamma p^2 &= Cx, \end{aligned}}
%CONT>>
